\section{Lecture 20 (March 3, 2025, est.): Frobenius type and lifting}
\label{lec20}
\noindent\textit{Scribe: Guillermo Sanmarco}

We continue over an algebraically closed field $k$ of characteristic $p>0$.
Our categories are semisimple unless otherwise specified. The Frobenius
functor takes values in $\mathcal C^{(1)}\boxtimes\mathsf{Ver}_p$, with the
scalar twist of Lecture~\ref{lec19}.

\subsection{Frobenius on the Verlinde category}
Assume $p>2$, and put $\varepsilon=L_{p-1}$. Recall the fusion rules
\[
 L_r\otimes L_s=
 \bigoplus_{i=1}^{\min(r,s,p-r,p-s)}L_{|r-s|+2i-1},
 \qquad \varepsilon\otimes L_s=L_{p-s}.
\]
In particular $\varepsilon^{\otimes2}=\mathbf1$, and
$\langle\varepsilon\rangle\simeq\SVec$. The subcategory
$\mathsf{Ver}_p^+$ consists of sums of the $L_s$ with $s$ odd.

\begin{proposition}\label{lec20-verlinde-frobenius}
Using the prime-field identification
$\mathsf{Ver}_p^{(1)}\simeq\mathsf{Ver}_p$, one has
\[
 \mathrm{Fr}(L_s)=
 \begin{cases}
  \mathbf1\boxtimes L_s,&s\text{ odd},\\
  \varepsilon\boxtimes L_{p-s},&s\text{ even}.
 \end{cases}
\]
\end{proposition}
\begin{proof}
Set $x=[L_2]$ in $K_0(\mathsf{Ver}_p)$. The recurrence
$[L_{s+1}]=x[L_s]-[L_{s-1}]$, together with $[L_0]=[L_p]=0$, gives
\[
 x^p=-2[\varepsilon]
 \quad\text{in }K_0(\mathsf{Ver}_p)\otimes k.
\]
One way to verify this polynomial identity is to put
$[L_s]=U_{s-1}(x/2)$, where
$U_0=1$, $U_1=2t$, and $U_{m+1}=2tU_m-U_{m-1}$, and reduce modulo
$p$ and $U_{p-1}(x/2)$.

Now $\operatorname{FPdim}(L_2)=2\cos(\pi/p)<2$. Preservation of this
dimension implies that $\mathrm{Fr}(L_2)$ is simple: every nonzero simple
summand has Frobenius--Perron dimension at least one. Write it as
$L_a\boxtimes L_b$. The identity of
Lemma~\ref{lec19-grothendieck-frobenius} says
$b[L_a]=-2[L_{p-1}]$ in $K_0(\mathsf{Ver}_p)\otimes k$.
Since $1\leq b<p$, comparison of basis coefficients gives
$a=p-1$ and $b=p-2$.
Finally apply $\mathrm{Fr}$ to
$L_2\otimes L_s=L_{s-1}\oplus L_{s+1}$ and use induction on $s$.
The displayed fusion rules verify the two parity cases and give the formula.
\end{proof}

\subsection{Frobenius type}
\begin{definition}
The \emph{Frobenius type} of $\mathcal C$ is the smallest fusion subcategory
$\mathcal A\subseteq\mathsf{Ver}_p$ such that
\[
 \mathrm{Fr}(\mathcal C)\subseteq\mathcal C^{(1)}\boxtimes\mathcal A.
\]
Equivalently, it is generated by all $L_r$ occurring as a second factor in
some simple summand of some $\mathrm{Fr}(X)$.
\end{definition}
For $p>3$, the possibilities are
$\Vect$, $\SVec$, $\mathsf{Ver}_p^+$, and $\mathsf{Ver}_p$.
To see why this list is exhaustive, a subcategory containing any
$L_r$ other than $\mathbf1$ or $\varepsilon$ contains $L_3$ as a summand of
$L_r\otimes L_r^*$. The fusion rules show that $L_3$ generates all the
odd-indexed simples. Adding any even-indexed simple then generates
$\varepsilon$ and hence all of $\mathsf{Ver}_p$.

\begin{example}
Both $\Vect$ and $\SVec$ have Frobenius type $\Vect$; for the odd line
$\pi$ and odd $p$, the $p$-cycle acts trivially on $\pi^{\otimes p}$.
Proposition~\ref{lec20-verlinde-frobenius} shows that
$\mathsf{Ver}_p$ has type $\mathsf{Ver}_p^+$. In particular the rank-two
category $\mathsf{Ver}_5^+$ has type $\mathsf{Ver}_5^+$, as computed in the
previous lecture. When $p=2$, $\mathsf{Ver}_2=\Vect$; when $p=3$,
$\mathsf{Ver}_3=\SVec$ and $\mathsf{Ver}_3^+=\Vect$.
\end{example}

\begin{proposition}\label{lec20-frobenius-naturality}
For an exact symmetric tensor functor $F:\mathcal C\to\mathcal D$, there
is a natural monoidal isomorphism
\[
 \mathrm{Fr}_{\mathcal D}\circ F
 \simeq(F^{(1)}\boxtimes\mathrm{id})\circ\mathrm{Fr}_{\mathcal C}.
\]
Consequently the Frobenius type of $\mathcal C$ is contained in that of
$\mathcal D$.
\end{proposition}
\begin{proof}
The symmetric structure of $F$ identifies the cyclic actions on
$F(X^{\otimes p})$ and $F(X)^{\otimes p}$. In the semisimple construction,
this identification is compatible with taking the multiplicity spaces and
with the quotient of $\Rep_k(C_p)$ by negligible maps. It gives the stated
isomorphism. Moreover an exact tensor functor is faithful, so it sends no
nonzero object to zero. Thus every second-factor summand occurring before
applying $F$ still occurs afterwards; it must belong to the type of
$\mathcal D$.
\end{proof}
In particular, a category admitting a super fiber functor has type
$\Vect$. If $\mathcal C$ itself has type $\Vect$, its Frobenius functor
lands in $\mathcal C^{(1)}$ and
\begin{equation}\label{lec20-type-vec-power}
 [\mathrm{Fr}(X)]=[X]^p
 \quad\text{in }K_0(\mathcal C)\otimes_{\mathbb Z}k.
\end{equation}

\subsection{Scalar extension and split categories}
For the discussion of lifting, temporarily let $F$ be an arbitrary field.
A finite semisimple $F$-linear category is \emph{split} if
$\End(S)=F$ for every simple object $S$. Equivalently, every simple remains
simple after extension to an algebraic closure. Scalar extension is formed
by extending morphism spaces and taking the additive idempotent completion;
we write $\mathcal C\otimes_F F'$ for a field extension $F'/F$.

For example, the real two-dimensional representation of $C_3$ in which a
generator acts by rotation through $2\pi/3$ is simple over $\mathbb R$.
Its endomorphism algebra is $\mathbb C$, and after extension to $\mathbb C$
it splits into two one-dimensional characters. Thus $\Rep_{\mathbb R}(C_3)$
is not split.

\subsection{The lifting theorem}
For a symmetric fusion category, define its \emph{global dimension} by
\[
 \dim(\mathcal C)=\sum_{S\text{ simple}}\dim(S)^2\in k.
\]
We call $\mathcal C$ \emph{nondegenerate} when this scalar is nonzero.
Here this word concerns global dimension; it does not mean that the
braiding has trivial symmetric center.

Return to an algebraically closed field $k$ of characteristic $p$, which
is in particular perfect, and let $W(k)$ be its ring of
Witt vectors. This is a complete discrete valuation ring of characteristic
zero with residue field $k$. A split lifting of a split fusion category
$\mathcal C$ to $W(k)$ consists of a rigid monoidal $W(k)$-linear category
with finite free morphism modules, the same simple labels and fusion
multiplicities, and reduction modulo $p$ equivalent to $\mathcal C$.
Extending to the fraction field gives a characteristic-zero fusion category.

\begin{theorem}[Etingof--Nikshych--Ostrik]
\label{lec20-lifting}
A nondegenerate fusion category over an algebraically closed field $k$ lifts to
$W(k)$, uniquely up to equivalence with prescribed reduction. A braiding,
and in particular a symmetry, lifts as well. Lifting is faithful on
equivalence classes of nondegenerate categories.
\end{theorem}
These are the lifting and faithfulness results of \cite[Section 9]{ENO}.
They are input theorems here; their proof uses vanishing of deformation
cohomology. The completeness and residue-field conditions on the base ring
are essential to this formulation: the theorem is not an assertion about
every local ring with residue field $k$.
